\scanpage{97}
\mspage{103}

En effet, on a pour $F$ loc. constant sur $X$

\[
\begin{aligned}
H^{\unsure{\cdot|i}}_{\mathbb D}(X,F)&\simeq H^{\cdot}_{\mathbb D}(Y,\underbrace{f_*(F)}_{\substack{\text{loc}\\\text{const. sur}\\Y}})
\xrightarrow[\substack{|\\ p\ \mathbb D\text{-équ.}\\\text{coh}}]{\sim}
H^{\cdot}_{\mathbb D}(Y',\underbrace{\unsure{p'|p^*}f_*(F)}_{\substack{\rotatebox{90}{$\simeq$}\\ f'_*q^{\unsure{*|\prime}}(F)\\\text{par prop. 2}}})\\
&\simeq H^{\cdot}_{\mathbb D}(Y',f'_*q^*F)\simeq H^{\cdot}_{\mathbb D}(X',q^*F)
\end{aligned}
\]

qed.

\uline{Prop 3} La notion de $\mathbb D$-fibration cohomologique \inserted{complète} est stable par chgt de base (trivial) et par composition.

\deleted{Il suffit de prouver les conditions a), b) pour $gf$ dans $X\xrightarrow{f}Y\xrightarrow{g}Z$.} Pour la stabilité par composition, il suffit de \deleted{\unsure{voir}} noter le cor. suivant de prop. 2

\uline{Cor. 2} (de Prop. 2) Soit $f:X\to Y$. Pour que $f$ soit une $\mathbb D$ fibration cohomologique \inserted{complète}, il f. et s. que pour $Y''\xrightarrow{\unsure{p'|p}}Y'\to Y$ \unsure{sur $Y$}, avec $p$ une $\mathbb D$-équiv coh, la flèche correspondante $X''=X\times_Y Y''\unsure{\relbar|\to}X'=X\times_Y\unsure{Y/|Y'}$ soit aussi une $\mathbb D$-équiv coh.

C'est nécessaire par le cor. 1, compte tenu que $f':X'=X\times_Y Y'\unsure{\relbar|\to}Y'$ \deleted{est} aussi une $\mathbb D$-fibr. coh. C'est suffisant, car il suffit de l'appliquer aux cas $\{y\}\unsure{\relbar|\to}Y_{/\{y\}}$, $\deleted{\unsure{\Rightarrow|y}}Y_{/\,\unsure{,y|y}}\unsure{\relbar|\to}Y_{/y'}$, qui sont des $\mathbb D'$-équiv. coh. en

% archive sheet number: 96
